
\subsection{Résumé}

\paragraph*{Liste des cours}\mbox{}\\[5pt]
{
\begin{tabular}{llccc}
      \hline
      \multicolumn{5}{c}{Vacations Université Paris-Saclay}                                                  \\
      \hline
      Année & Intitulé                                             & TD & TP  & HETD\footnote{\label{TPTDUL} %
            L'Université Paris-Saclay ne fait généralement pas la différence entre TP et TD.
            Je fais la distinction quand elle est possible mais je garde la comptabilisation
      des HETD.}                                                                                             \\
      \hline
      \shortstack[l]{                                                                                        \\ 2021-2022}
            & L3 Programmation parallèle (C++/OpenMP)              &    & 2h  & 2h                           \\
      \hline
      \multirow{3}{*}{\shortstack[l]{                                                                        \\ 2022-2023}}
            & QDCS Calcul haute performance (C++/Cuda)             &    & 11h & 11h                          \\
            & IoT Programmation MPI  (C++/MPI)                     &    & 14h & 14h                          \\
            & L3 Programmation parallèle (C++/OpenMP)              & 2h &     & 2h                           \\
      \hline
      \multirow{2}{*}{\shortstack[l]{                                                                        \\ 2023-2024}}
            & Polytech - Programmation parallèle (C++/OpenMP/Cuda) &    & 20h & 20h                          \\
            & IoT Programmation MPI (C++/MPI)                      &    & 12h & 12h                          \\
            & L3 Programmation parallèle (C++/OpenMP)              &    & 6h  & 6h                           \\
      \hline
      \multirow{2}{*}{\shortstack[l]{                                                                        \\ 2025-2026}}
            & L3 Calcul Scientifique (Matlab)                      &    & 10h & 10h                          \\
            & L3 Projet Calcul Scientifique (Julia)                &    & 12h & 12h                          \\
      \hline
      \multicolumn{5}{c}{\bfseries TOTAL : 2h TD - 87h TP
      $\Longrightarrow$ 89h HETD\footnoteref{TPTDUL} }                                                       \\
      \hline
\end{tabular}%
}

\paragraph*{Responsabilités}\mbox{}\\[44pt]
\begin{center}
      \begin{tabular}{|l|c|c|c|c|c|c|c|}
            \cline{1-1}
            \multicolumn{1}{|c}{Module}                                                &
            \multicolumn{1}{|c}{\begin{rotate}{45} Responsable du module \end{rotate}} &
            \multicolumn{1}{c}{\begin{rotate}{45} Enseignement \end{rotate}}           &
            \multicolumn{1}{c}{\begin{rotate}{45} Conception d'exercice \end{rotate}}  &
            \multicolumn{1}{c}{\begin{rotate}{45} Encadrement de TD/TP \end{rotate}}   &
            \multicolumn{1}{c}{\begin{rotate}{45} Conception d'examen \end{rotate}}    &
            \multicolumn{1}{c}{\begin{rotate}{45} Surveillance d'examen \end{rotate}}  &
            \multicolumn{1}{c}{\begin{rotate}{45} Corrections \end{rotate}}                                                                                                \\
            \hline
            L3 Programmation parallèle                                                 &           &           &           & \ding{51} &           &           &           \\
            \hline
            \hline
            QDCS Calcul haute performance                                              &           &           &           & \ding{51} &           &           & \ding{51} \\
            \hline
            IoT Programmation MPI                                                      &           & \ding{51} & \ding{51} & \ding{51} &           & \ding{51} & \ding{51} \\
            \hline
            L3 Programmation parallèle                                                 &           &           &           & \ding{51} &           & \ding{51} & \ding{51} \\
            \hline
            Polytech - Algèbre Linéaire                                                &           &           &           &           &           & \ding{51} &           \\
            \hline
            \hline
            Polytech - Programmation parallèle                                         &           & \ding{51} &           & \ding{51} &           & \ding{51} & \ding{51} \\
            \hline
            IoT Programmation MPI                                                      &           & \ding{51} & \ding{51} & \ding{51} &           & \ding{51} & \ding{51} \\
            \hline
            L3 Programmation parallèle                                                 &           &           &           & \ding{51} &           & \ding{51} & \ding{51} \\
            \hline
            \hline
            Calcul Scientifique                                                        &           &           &           & \ding{51} & \ding{51} & \ding{51} & \ding{51} \\
            \hline
            Projet Calcul Scientifique                                                 & \ding{51} & \ding{51} & \ding{51} & \ding{51} & \ding{51} & \ding{51} & \ding{51} \\
            \hline
      \end{tabular}%
\end{center}
\mbox{}\\

\subsection{Description des Vacations}

\small

\paragraph*{2021 -- 2022}
\begin{description}
      \item[Intitulé :] \textit{\textbf{Programmation parallèle}}
      \item[Responsable :] Oguz Kaya
            (\href{mailto:oguz.kaya@universite-paris-saclay.fr}{oguz.kaya@universite-paris-saclay.fr})
      \item[Public :] Troisième année de Licence d'informatique, UPSaclay Orsay (1 Groupe
            de TP : $\sim 20$ étudiants)
      \item[Volume horaire :] 2h TP
      \item[Contenu :] Notions classiques d'algorithmique parallèle (OpenMP, mémoire
            partagé), programmation en C++
      \item[Responsabilités :] Encadrement de TP
\end{description}

\paragraph*{2022 -- 2023}
\begin{description}
      \item[Intitulé :] \textit{\textbf{Calcul haute performance}}
      \item[Responsable :] Oguz Kaya
            (\href{mailto:oguz.kaya@universite-paris-saclay.fr}{oguz.kaya@universite-paris-saclay.fr})
      \item[Public :] Première année de Master d'informatique QDCS, UPSaclay Orsay (1
            Groupe de TP : $\sim 20$ étudiants)
      \item[Volume horaire :] 11h TP
      \item[Contenu :] Notions d'utilisation de cluster et d'algorithmes parallèles
            optimisés pour l'architecture (en C++).
      \item[Responsabilités :] Encadrement de TP, corrections de l'examen.
\end{description}

\begin{description}
      \item[Intitulé :] \textit{\textbf{Programation MPI}}
      \item[Responsable :] Karim Hasnaoui
            (\href{mailto:karim.hasnaoui@ijclab.in2p3.fr}{karim.hasnaoui@ijclab.in2p3.fr})
      \item[Public :] Première année de Master d'informatique IoT, UPSaclay Orsay (1 Groupe
            de TP : $\sim 20$ étudiants)
      \item[Volume horaire :] 14h TP
      \item[Contenu :] Notions d'utilisation d'algorithme distribué et de la mémoire
            distribué. (en C++ avec la bibliothèque MPI). Utilisation du cluster Ruche pour
            application.
      \item[Responsabilités :] Encadrement du cours ainsi que tu TD/TP. Conception
            d'exercice de TP ainsi que surveillance et correction de l'examen.
\end{description}

\begin{description}
      \item[Intitulé :] \textit{\textbf{Programmation parallèle}}
      \item[Responsable :] Oguz Kaya
            (\href{mailto:oguz.kaya@universite-paris-saclay.fr}{oguz.kaya@universite-paris-saclay.fr})
      \item[Public :] Troisième année de Licence d'informatique, UPSaclay Orsay (1 Groupe
            de TP : $\sim 20$ étudiants)
      \item[Volume horaire :] 2h TD
      \item[Contenu :] Notions classiques d'algorithmique parallèle (OpenMP, mémoire
            partagé), programmation en C++
      \item[Responsabilités :] Encadrement de TD ainsi que surveillance et correction de
            l'examen.
\end{description}

\paragraph*{2023 -- 2024}
\begin{description}
      \item[Intitulé :] \textit{\textbf{Programmation parallèle}}
      \item[Responsable :] Marc Baboulin
            (\href{mailto:baboulin@lmf.cnrs.fr}{baboulin@lmf.cnrs.fr})
      \item[Public :] Deuxième année d'école d'ingénieur, Polytechnique Orsay (1 Groupe de
            TP : $\sim 20$ étudiants)
      \item[Volume horaire :] 20h TP
      \item[Contenu :] Notions classiques d'algorithmique parallèle (OpenMP, mémoire
            partagé), programmation en C++
      \item[Responsabilités :] Enseignement du cours. Et encadrement de TP ainsi que
            surveillance et correction de l'examen.
\end{description}

\begin{description}
      \item[Intitulé :] \textit{\textbf{Programation MPI}}
      \item[Responsable :] Karim Hasnaoui
            (\href{mailto:karim.hasnaoui@ijclab.in2p3.fr}{karim.hasnaoui@ijclab.in2p3.fr})
      \item[Public :] Première année de Master d'informatique IoT, UPSaclay Orsay (1 Groupe
            de TP : $\sim 20$ étudiants)
      \item[Volume horaire :] 12h TP
      \item[Contenu :] Notions d'utilisation d'algorithme distribué et de la mémoire
            distribué. (en C++ avec la bibliothèque MPI). Utilisation du cluster Ruche pour
            application.
      \item[Responsabilités :] Encadrement du cours ainsi que tu TD/TP. Conception
            d'exercice de TP ainsi que surveillance et correction de l'examen.
\end{description}

\begin{description}
      \item[Intitulé :] \textit{\textbf{Programmation parallèle}}
      \item[Responsable :] Oguz Kaya
            (\href{mailto:oguz.kaya@universite-paris-saclay.fr}{oguz.kaya@universite-paris-saclay.fr})
      \item[Public :] Troisième année de Licence d'informatique, UPSaclay Orsay (1 Groupe
            de TP : $\sim 20$ étudiants)
      \item[Volume horaire :] 6h TP
      \item[Contenu :] Notions classiques d'algorithmique parallèle (OpenMP, mémoire
            partagé), programmation en C++
      \item[Responsabilités :] Encadrement de TD ainsi que surveillance et correction de
            l'examen.
\end{description}

\paragraph*{2025 -- 2026}
\begin{description}
      \item[Intitulé :] \textit{\textbf{Calcul Scientifique}}
      \item[Responsable :] Ehouarn Simon
            (\href{mailto:ehouarn.simon@irit.fr}{ehouarn.simon@irit.fr})
      \item[Public :] Première année d'école d'ingénieur, Enseeith Toulouse (3 Groupe de TP
            : $\sim 45$ étudiants)
      \item[Volume horaire :] 10h TP
      \item[Contenu :] Notions classiques d'algèbre linéaire (Gram-Schmith, factorisation
            LU), programmation en Matlab
      \item[Responsabilités :] Encadrement de TP.
\end{description}

\begin{description}
      \item[Intitulé :] \textit{\textbf{Projet Calcul Scientifique}}
      \item[Responsable :] \underline{\textbf{Matthieu Robeyns}}
            (\href{mailto:matthieu.robeyns@irit.fr}{matthieu.robeyns@irit.fr})
      \item[Public :] Première année d'école d'ingénieur, Enseeith Toulouse (4 Groupe de TP
            : $\sim 80$ étudiants)
      \item[Volume horaire :] 7h TP
      \item[Contenu :] Notions d'utilisation d'algorithme randomisé pour des matrices et
            découverte des tenseurs avec l'approximation de rang faibles (en Julia avec
            l'interface Pluto).
      \item[Responsabilités :] Création du projet et de ses questions. Encadrement du
            cours/TP.
\end{description}

\subsection{Responsable de Stage}

\paragraph*{Mai -- Juillet 2023 :}
Supervision d'un stagiaire de M1 sur l'intégration de la bibliothèque MPLAPACK dans Celeste pour simuler les algorithmes de mixed precision.
J'ai appris à encadrer et écouter un stagiaire pendant 6 mois pour avancer ensemble dans la continuer de ma thèse.
Nous avons ensemble réussi à joindre l'aspect précision mixte de MPLAPACK via une bibliothèque MPREAL.
Puis nous l'avons implémenter dans un dossier de Celeste qui à ce jour n'est pas publique.
\normalsize

\subsection{Responsable d'un séminaire de groupe SéPAG}

\paragraph*{Septembre 2022 -- Decembre 2024 ($\simeq$60h) :}
Création plus organisation du Séminaire SéPAG réunissant trois équipes (A\&O, Galac, \underline{ParSys}) du groupa AAC du laboratoire LISN à Paris-Saclay.
Avec l'aide de Nathan SOLAL (A\&O) et Loïc LE MOGNE (Galac), nous avons organisé des réunions mensuelles permettant à des membres de l'équipes de communiquer leurs découvertes au sein des autres équipes.
J'ai été le responsable de l'équipe Parsys dans ce projet pendant deux années consécutives et j'ai présenté trois travaux de recherches.\\
Lien du site : \href{https://hebergement.universite-paris-saclay.fr/sepag/}{SéPAG}
\normalsize

\subsection{Projet d'enseignement}
Au-delà de la recherche, un poste universitaire repose avant tout sur la
transmission des savoirs à travers l'enseignement, je dirais même qu'il
constitue la clé de voûte de l'activité de recherche. C'est en effet par
l'enseignement que nous partageons nos connaissances et notre passion pour nos
thématiques scientifiques, suscitant à leur tour l'intérêt des étudiants pour
la recherche. Initialement, je souhaitais devenir professeur de mathématiques,
cependant, en suivant des cours d'informatique encadrés par une enseignante
passionnée, mon intérêt pour l'enseignement s'est élargi, m'amenant à vouloir
transmettre à la fois les mathématiques et l'informatique ; dans cet esprit de
résonance et de partage des connaissances entre l'étudiant et l'enseignant. Je
pense que l'enseignement permet également de revisiter les notions
fondamentales des disciplines que nous transmettons, mais aussi d'en acquérir
de nouvelles, notamment à travers l'encadrement de séances de cours ou TD, ou
les échanges avec les collègues. Selon moi, il contribue à renforcer la
cohésion des équipes et favorise l'émergence de projets pédagogiques
transversaux au sein d'un laboratoire.

L'enseignement que je voudrais donner au sein de l'équipe \equipe constituerait
une opportunité d'enrichir mes compétences en algèbre linéaire et en calcul
haute performance (HPC), tout en y apportant mon expertise. Je pourrais
notamment introduire des notions telles que le calcul tensoriel ou le calcul en
précision mixte appliqué aux approximations de rang faible. Par ailleurs, sous
réserve des besoins de l'équipe, je pourrais contribuer à des enseignements
existants en y intégrant des mises à jour, par exemple sur les architectures
récentes et la programmation sur GPU.

L'enseignement est aujourd'hui fortement impacté par l'IA, au point que
certaines méthodes traditionnelles peuvent paraître obsolètes. Pour maintenir
une approche stimulante pour les étudiants, je privilégierais un enseignement
en groupe, débutant par des outils simples ---un papier, un crayon et une bonne
motivation--- afin d'aborder la résolution de problèmes donnés. Une fois cette
phase de réflexion menée en interaction avec l'enseignant terminée, nous
pourrions passer à la partie pratique sur ordinateur. Pour vérifier la
compréhension des notions, des rendus de projets peuvent être remis, mais
accompagnés d'un contrôle écrit visant principalement à résumer le travail
réalisé. Des évaluations sous forme d'exposés pourraient également être
envisagées, dans un format proche de celui d'une conférence : 5 à 10 minutes de
présentation suivies de 5 minutes de questions, permettant de vérifier que les
concepts ont été assimilés et préparant les étudiants à un format plus proche
de leur avenir s'ils s'orientent dans le domaine de la recherche.

Le projet que j'ai créé ``Projet Calcul Scientifique'' sur les méthodes
randomisées pour les matrices et tenseurs, durant mon post-doctorat, m'a fait
prendre conscience que j'ai les capacités pour créer de nouvelles unités
d'enseignement, bien que je ne sois pas réticent à travailler sur les
différentes thématiques suivantes.

\begin{flushleft}
\textbf{Calcul haute performance} Mon master PDCS de l'Université Paris-Saclay m'a permis d'apprendre diverses méthodes de calcul haute performance et de les employer dans divers clusters. Avec mon expérience dans le domaine et sur les clusters tels que PlafRim, Ruche, Jean Zay, Jupyter, je pense pourvoir enseigner les rudiments de la programmation avec clusters (utilisation de slurm par exemple).
\end{flushleft}

\begin{flushleft}
\textbf{Programmation Parallèle} Depuis mon master j'ai compris que la programmation séquentielle est essentielle, mais pour obtenir de bonnes performances dans des domaines complexes tels que la physique, la parallélisation devient primordiale. Depuis cette période j'ai toujours gardé, de près ou de loin, une expérience avec le calcul parallèle, que ce soit avec OpenMP, MPI ou bien les graphes de tâches. Devenant un érudit du domaine, je peux facilement jongler entre l'aspect CPU et/ou GPU pour enseigner dans des systèmes hétérogènes ou homogènes.
\end{flushleft}

\begin{flushleft}
\textbf{Calcul Scientifique} Il est évident que coder sans savoir ce que le code va donner n'est pas pertinent. C'est pourquoi je pense que la matière calcul scientifique est un enseignement pilier d'un parcours informatique (surtout HPC) et qu'il ne faut pas la négliger. Prévoir le résultat, ou bien le prouver, ou à minima le comprendre est un apprentissage qui se fait dans le module "Calcul Scientifique" que je veux bien enseigner aux étudiants pour comprendre la liaison entre l'algèbre linéaire et l'informatique. On pourrait leur enseigner la base, avec la méthode de Gram-Schmidt, puis la factorisation LU et QR pour ensuite tester d'autres méthodes voir même des approximations de rang faible.
\end{flushleft}

\begin{flushleft}
\textbf{Algorithme en précision mixte} Durant ma thèse j'ai pu explorer un domaine qui n'était que purement théorique pendant mon parcours, c'est la précision mixte. Celle-ci m'a ouvert tellement de possibilité d'implémentation de code, inenvisageable en haute précision. Il est pour moi pertinent de l'enseigner dans des cursus HPC, surtout avec l'avenir des GPU qui tend de plus en plus vers la précision faible (voir les integer 8 !). J'ai acquéri ces connaissances avec des personnes expertes dans ce domaine et j'aimerais leur faire honneur en enseignant à mon tour ce qu'ils m'ont transmis. Il serait très agréable de jouer avec les différentes précisions standard du format IEEE tel que FP16, FP32, FP64, pour expliquer aux étudiants l'impact des calculs dans certaines précisions et les erreurs associées.
\end{flushleft}

\begin{flushleft}
\textbf{Complexité algorithmique} Les calculs de complexité ne sont pas très représentés dans les parcours pour les étudiants. Pourtant ils permettent très bien d'expliquer que deux algorithmes donnant le même résultat peuvent avoir des temps d'exécution totalement différents. Mon parcours m'a rigoureusement impliqué dans le calcul de complexité pour pouvoir vérifier que les méthodes que j'implémentais était cohérentes et aussi pour pouvoir se comparer avec l'état de l'art. Je trouverais formateur d'enseigner cette notion aux étudiants au travers de calculs simples, et de finir le cours en autonomie pour faire une implémentation presque équivalente ou voir plus rapide à des méthodes étudiées durant leur parcours.
\end{flushleft}

Cette liste présente un résumé des différentes thématiques et des modules dans
lesquels je dispose des connaissances et de l'expérience pour prendre
rapidement des responsabilités dans les différentes formations proposées par le
département \departement. Je reste néanmoins parfaitement ouvert à
l'enseignement sur d'autres thématiques. D'autre part, j'aimerais mettre à
profit mon expérience professionnelle en m'engageant dans les activités
d'échange avec les entreprises, comme par exemple en participant aux activités
de suivi des étudiants en alternance.

% \begin{itemize}
% \item Calcul haute performance : Master 1 et 2
% \item Programmation Parallèle : Licence 3 ou Masters 1 et 2 ou Ingénieur
% \item Algèbre Linéaire : Licence 3 ou Masters 1 et 2 ou Ingénieur
% \item Algorithme de précision mixte : Master 1 et 2
% \item Complexité algorithmique : Licence 2 et 3
% \end{itemize}
